\section{Khe hở giữa trị riêng và giá trị kỳ dị}
\label{sec:ch03-khe-ho}

\index{bán kính phổ!so với chuẩn phổ}%
Cho tới đây tôi đã viết điều kiện qua \(\norm{\mat{W}_{hh}}\), \tn{chuẩn phổ}{spectral norm}. Bài
báo không viết thế. Bài phát biểu \tn{điều kiện đủ}{sufficient condition} qua \(\lambda_1\), mà họ định
nghĩa bằng chữ là \enquote{trị tuyệt đối của \tn{trị riêng}{eigenvalue} lớn nhất} của
\(\mat{W}_{hh}\), tức \tn{bán kính phổ}{spectral radius} \(\rho(\mat{W}_{hh})\).

Hai số đó không bằng nhau. Quan hệ giữa chúng chỉ có một chiều:

\begin{equation}
  \rho(\mat{W}) \le \norm{\mat{W}},
  \label{eq:ch03-ro-nho-hon-chuan}
\end{equation}

và dấu bằng xảy ra khi \(\mat{W}\) là \tn{ma trận chuẩn tắc}{normal matrix}, tức
\(\mat{W}\mat{W}\transpose = \mat{W}\transpose\mat{W}\). Ma trận trực giao là
chuẩn tắc, ma trận đối xứng là chuẩn tắc. Một ma trận trọng số hồi quy được
khởi tạo ngẫu nhiên rồi huấn luyện thì không.

Chỗ lệch nằm ở đây. Chứng minh của bài, tức bất đẳng
thức~\eqref{eq:ch03-chan-tung-thua-so}, chặn \(\norm{\mat{W}_{hh}}\): đó là một
\tn{giá trị kỳ dị}{singular value}. Phát biểu bằng lời của bài nói về trị riêng.
Với ma trận chuẩn tắc hai cách đọc trùng nhau và không ai nhận ra. Với ma trận
không chuẩn tắc, \(\rho(\mat{W}_{hh}) < 1/\gamma\) \emph{không} kéo
theo~\eqref{eq:ch03-chan-tung-thua-so}, nên chứng minh không chạy.

Câu hỏi là chỗ lệch ấy có hậu quả thật hay chỉ là một dấu sao trong cước chú.
Tôi dựng ma trận rẻ nhất tách được hai số đó:

\begin{minted}{python}
def jordan_block(n: int, eigenvalue: float, off_diagonal: float) -> torch.Tensor:
    w = torch.eye(n, dtype=torch.float64) * eigenvalue
    idx = torch.arange(n - 1)
    w[idx, idx + 1] = off_diagonal
    return w
\end{minted}

\index{ma trận chuẩn tắc!phản ví dụ}%
Ma trận tam giác trên với \(0.9\) trên đường chéo và \(3.0\) ngay phía trên nó.
Mọi trị riêng đều bằng \(0.9\), nên bán kính phổ là \(0.9\) bất kể phần tử ngoài
đường chéo lớn đến đâu. Chuẩn phổ thì lớn theo. Để tách phần đại số tuyến tính
khỏi bão hoà, tôi chạy nó ở trường hợp tuyến tính, tức \(\sigma\) là ánh xạ đồng
nhất, đúng bối cảnh phần phụ lục của bài: khi đó tích
trong~\eqref{eq:ch03-tich-jacobian} thu về đúng luỹ thừa ma trận
\(\mat{W}_{hh}^{\,l}\) với \(l = t-k\).

\begin{measured}{tích Jacobian tăng dưới một bán kính phổ nhỏ hơn 1}
\begin{minted}[fontsize=\footnotesize]{text}
W = [[0.9, 3.0], [0, 0.9]]
spectral radius (largest |eigenvalue|) = 0.900000
spectral norm   (largest singular value) = 3.249286

l     ||W^l||
1     3.249286e+00
5     9.876803e+00
9     1.163551e+01
10    1.163307e+01
20    8.106934e+00
50    8.589935e-01
100   8.853879e-03

peak at l=9, ||W^l||=11.635514
amplification over l=1: 3.5809
first l with ||W^l|| < 1: 49
\end{minted}

Đo bằng \code{experiments/ch03\_nonnormal.py} tại tag \code{ch03}.
\end{measured}

Bán kính phổ là \(0.9\), tức điều kiện đủ của bài, đọc theo đúng chữ họ viết,
được thoả mãn. Vậy mà tích Jacobian \emph{tăng} suốt chín bước, đạt 11.635514,
gấp 3.5809 lần giá trị của chính nó ở bước 1, và phải tới bước 49 mới lần đầu
tụt xuống dưới 1.

Kiểm lại bằng tay được, và nên kiểm: phần tử ngoài đường chéo của
\(\mat{W}^{\,l}\) là \(l \cdot 3 \cdot 0.9^{\,l-1}\), đạo hàm theo \(l\) triệt
tiêu tại \(l = -1/\ln 0.9\), tức khoảng 9.49. Đỉnh rơi vào giữa bước 9 và bước
10, và số đo cho hai giá trị gần bằng nhau ở đúng hai bước đó.

\subsection*{Phát biểu đúng là gì}

\index{vanishing gradient!phát biểu chính xác}%
Không nên đọc chỗ này thành \enquote{bài báo sai}. Trong trường hợp tuyến tính,
đúng trường hợp vừa đo, phát biểu của họ đúng ở nghĩa tiệm cận: công thức
Gelfand cho \(\norm{\mat{W}^{\,l}}^{1/l} \to \rho(\mat{W})\), nên
\(\rho(\mat{W}_{hh}) < 1\) vẫn kéo theo \(\norm{\mat{W}_{hh}^{\,l}} \to 0\) và
\tn{độ dốc}{gradient} rồi sẽ tắt. Cái không đúng là chặn đơn điệu \(\eta^{\,t-k}\)
trong~\eqref{eq:ch03-chan-mu}. Chặn đó cần giá trị kỳ dị, và chỉ có nó mới cho
phép nói gradient tắt \emph{ngay từ bước đầu}.

Với \(\sigma\) phi tuyến thì ngay cả phần tiệm cận cũng không suy ra được theo
kiểu ấy, vì tích trong~\eqref{eq:ch03-tich-jacobian} là tích của những ma trận
\emph{khác nhau}, mỗi bước một \(\diag(\sigma'(\vect{a}_{i-1}))\) riêng, chứ
không phải luỹ thừa của một ma trận. Công thức Gelfand không nói gì về tích như
thế.

Khoảng cách giữa hai phát biểu là khoảng cách giữa \enquote{tắt} và
\enquote{tắt sau bao lâu}, và với chuỗi hữu hạn thì đó là toàn bộ vấn đề. Tích
trên tăng tới bước 9, rồi giảm, nhưng phải tới bước 49 mới lần đầu xuống dưới 1.
Trên một chuỗi dài 50 bước, không có bước nào mà tín hiệu lỗi đi ngược bị co
lại: nó bị khuếch đại ở mọi khoảng cách mà chuỗi ấy có. Đó không phải vanishing
gradient, dưới một bán kính phổ mà lý thuyết đọc lướt bảo là an toàn.

Đây là lý do sách tách \(\rho(\mat{W}_{hh})\) và \(\norm{\mat{W}_{hh}}\) thành
hai ký hiệu và không bao giờ dùng lẫn. Khi bạn đọc một bài báo nói
\enquote{spectral radius nhỏ hơn 1 nên gradient tắt}, câu hỏi cần đặt là ma
trận có chuẩn tắc không, và câu trả lời gần như luôn là không.
